% Options for packages loaded elsewhere
\PassOptionsToPackage{unicode}{hyperref}
\PassOptionsToPackage{hyphens}{url}
\PassOptionsToPackage{dvipsnames,svgnames,x11names}{xcolor}
\documentclass[
  11pt,
]{article}
\usepackage{xcolor}
\usepackage[margin=25mm]{geometry}
\usepackage{amsmath,amssymb}
\setcounter{secnumdepth}{-\maxdimen} % remove section numbering
\usepackage{iftex}
\ifPDFTeX
  \usepackage[T1]{fontenc}
  \usepackage[utf8]{inputenc}
  \usepackage{textcomp} % provide euro and other symbols
\else % if luatex or xetex
  \usepackage{unicode-math} % this also loads fontspec
  \defaultfontfeatures{Scale=MatchLowercase}
  \defaultfontfeatures[\rmfamily]{Ligatures=TeX,Scale=1}
\fi
\usepackage{lmodern}
\ifPDFTeX\else
  % xetex/luatex font selection
\fi
% Use upquote if available, for straight quotes in verbatim environments
\IfFileExists{upquote.sty}{\usepackage{upquote}}{}
\IfFileExists{microtype.sty}{% use microtype if available
  \usepackage[]{microtype}
  \UseMicrotypeSet[protrusion]{basicmath} % disable protrusion for tt fonts
}{}
\makeatletter
\@ifundefined{KOMAClassName}{% if non-KOMA class
  \IfFileExists{parskip.sty}{%
    \usepackage{parskip}
  }{% else
    \setlength{\parindent}{0pt}
    \setlength{\parskip}{6pt plus 2pt minus 1pt}}
}{% if KOMA class
  \KOMAoptions{parskip=half}}
\makeatother
\usepackage{listings}
\usepackage{seqsplit}
\usepackage{needspace}
\newcommand{\passthrough}[1]{#1}
\lstset{defaultdialect=[5.3]Lua}
\lstset{defaultdialect=[x86masm]Assembler}
\usepackage{longtable,booktabs,array}
\usepackage{calc} % for calculating minipage widths
% Correct order of tables after \paragraph or \subparagraph
\usepackage{etoolbox}
\AtBeginEnvironment{longtable}{\small}
\makeatletter
\patchcmd\longtable{\par}{\if@noskipsec\mbox{}\fi\par}{}{}
\makeatother
% Allow footnotes in longtable head/foot
\IfFileExists{footnotehyper.sty}{\usepackage{footnotehyper}}{\usepackage{footnote}}
\makesavenoteenv{longtable}
\setlength{\emergencystretch}{3em} % prevent overfull lines
\providecommand{\tightlist}{%
  \setlength{\itemsep}{0pt}\setlength{\parskip}{0pt}}
\emergencystretch=3em
\AtBeginDocument{\Urlmuskip=0mu plus 2mu}
\lstset{breaklines=true,basicstyle=\small\ttfamily,columns=fullflexible,keepspaces=true}
\setlength{\tabcolsep}{4pt}
\usepackage{iftex}
\ifXeTeX
  \usepackage{newunicodechar}
  \makeatletter
  \newcommand{\DeclareUnicodeCharacter}[2]{%
    \begingroup
    \uccode`\~="#1\relax
    \uppercase{\def\temp{~}}%
    \expandafter\endgroup
    \expandafter\newunicodechar\expandafter{\temp}{#2}%
  }
  \makeatother
\fi
\ifdefined\DeclareUnicodeCharacter
\DeclareUnicodeCharacter{2192}{\ensuremath{\to}}
\DeclareUnicodeCharacter{2190}{\ensuremath{\leftarrow}}
\DeclareUnicodeCharacter{2194}{\ensuremath{\leftrightarrow}}
\DeclareUnicodeCharacter{21D2}{\ensuremath{\Rightarrow}}
\DeclareUnicodeCharacter{21A6}{\ensuremath{\mapsto}}
\DeclareUnicodeCharacter{2212}{\ensuremath{-}}
\DeclareUnicodeCharacter{2264}{\ensuremath{\leq}}
\DeclareUnicodeCharacter{2265}{\ensuremath{\geq}}
\DeclareUnicodeCharacter{2260}{\ensuremath{\neq}}
\DeclareUnicodeCharacter{2113}{\ensuremath{\ell}}
\DeclareUnicodeCharacter{03A3}{\ensuremath{\Sigma}}
\DeclareUnicodeCharacter{03A6}{\ensuremath{\Phi}}
\DeclareUnicodeCharacter{0393}{\ensuremath{\Gamma}}
\DeclareUnicodeCharacter{221E}{\ensuremath{\infty}}
\DeclareUnicodeCharacter{2208}{\ensuremath{\in}}
\DeclareUnicodeCharacter{2209}{\ensuremath{\notin}}
\DeclareUnicodeCharacter{211A}{\ensuremath{\mathbb{Q}}}
\DeclareUnicodeCharacter{211D}{\ensuremath{\mathbb{R}}}
\DeclareUnicodeCharacter{2124}{\ensuremath{\mathbb{Z}}}
\DeclareUnicodeCharacter{2115}{\ensuremath{\mathbb{N}}}
\DeclareUnicodeCharacter{2200}{\ensuremath{\forall}}
\DeclareUnicodeCharacter{2203}{\ensuremath{\exists}}
\DeclareUnicodeCharacter{2286}{\ensuremath{\subseteq}}
\DeclareUnicodeCharacter{00D7}{\ensuremath{\times}}
\DeclareUnicodeCharacter{03B4}{\ensuremath{\delta}}
\DeclareUnicodeCharacter{2202}{\ensuremath{\partial}}
\DeclareUnicodeCharacter{03B2}{\ensuremath{\beta}}
\DeclareUnicodeCharacter{03B1}{\ensuremath{\alpha}}
\DeclareUnicodeCharacter{03C6}{\ensuremath{\phi}}
\DeclareUnicodeCharacter{03B3}{\ensuremath{\gamma}}
\DeclareUnicodeCharacter{00B9}{\ensuremath{^1}}
\DeclareUnicodeCharacter{00B2}{\ensuremath{^2}}
\DeclareUnicodeCharacter{2070}{\ensuremath{^0}}
\DeclareUnicodeCharacter{2080}{\ensuremath{_0}}
\DeclareUnicodeCharacter{2081}{\ensuremath{_1}}
\DeclareUnicodeCharacter{2082}{\ensuremath{_2}}
\fi
\usepackage{bookmark}
\IfFileExists{xurl.sty}{\usepackage{xurl}}{} % add URL line breaks if available
\urlstyle{same}
\hypersetup{
  pdftitle={PFP X: Exact Host Checking of Model Verdicts on Small Finite Algebraic Challenges --- Lattice Holonomy and Six Formal-Mathematics Pilots},
  pdfauthor={Jeremy H. Carroll},
  colorlinks=true,
  linkcolor={Maroon},
  filecolor={Maroon},
  citecolor={Blue},
  urlcolor={Blue},
  pdfcreator={LaTeX via pandoc}}

\title{PFP X: Exact Host Checking of Model Verdicts on Small Finite Algebraic Challenges — Lattice Holonomy and Six Formal-Mathematics Pilots}
\author{Jeremy H. Carroll}
\date{2026-10-02 · Version 0.1.0 · Research methods paper · Author-review revision}
\begin{document}
\makeatletter
\begin{center}
{\LARGE \@title\par}
\vspace{0.8em}
{\normalsize \@author\par}
\vspace{0.5em}
{\small \@date\par}
\end{center}
\makeatother
Code baseline:
\texttt{\seqsplit{a623568eb0f853b968bc419bf06dfc6af256ab02}}. Each model
run discussed in Sections 3 and 4 is cited by its repository path and by
the commit that recorded it; all of those commits precede this baseline.

Version 0.1.0 designates the initial public preprint candidate. Earlier
1.x and 2.x numbers in the change log identify internal drafts and
corrections, not prior Zenodo publications. No DOI or external
publication is claimed.

\hypertarget{abstract}{%
\subsection{Abstract}\label{abstract}}

Exact host checking separates a model's proposed verdict from a verified
certificate. For finite-group connections on square lattices, this paper
provides a complete Python flatness checker over a validated
multiplication table. It checks ordered plaquette products with inverses
on reversed edges; on a torus, two cycle holonomies distinguish
gauge-trivial from other flat connections. Skipping four-identity-link
plaquettes is an optimization whose saving depends on the gauge. On five
retained C017 cases, fresh host replay reproduces the expected verdicts
and the historical harness decisions on all 15 retained solution
records: 13 parsed certificates and two parser-error records, with 10
accepted and 5 rejected overall. Ten further historical responses retain
only labels and reasons. The five runs accepted 0, 0, 3, 3 and 4
responses out of five under different models, prompts and budgets; they
do not identify a causal prompt effect. Six formal-mathematics pilots
report three models on ten cases per cell, with two C049 panels and two
C050 runs. The author's case-by-case annotation identifies
answer-bearing identifiers in 36 of 60 reconstructable panel cases. In
two challenges a constant refusal passes all ten under the original
checkers. These counts therefore do not measure unaided performance. The
package supplies retained reports, offline table recounts, negative
controls, and one narrowly scoped Lean identity-link lemma. Unsupported
benchmark and effect claims from the initial internal draft are
withdrawn. No model run is repeated for this revision.

\hypertarget{problem-and-design}{%
\subsection{1. Problem and design}\label{problem-and-design}}

The Proposal Fidelity Protocol (PFP), developed in this collection under
the legacy name Symbol Game, separates a candidate answer from the host
check that decides whether the answer is accepted. This paper applies
that separation to model answers about small finite algebraic objects.
The model states a verdict and a certificate. A deterministic host
program recomputes the property from the problem data using only exact
operations: table lookups and integer arithmetic. The decision depends
on the certificate, not on the model's explanation.

Lattice holonomy is a useful test of this design. For a nonabelian
group, the product around a plaquette depends on the order of
multiplication and on the orientation, and an error in either gives a
confident wrong answer. Checking is cheap: one pass over the plaquettes
and, on a torus, two cycle products. The task therefore shows both
whether a host catches order errors and what a harness must tell a model
for its answer to be checkable at all.

The paper makes three contributions:

\begin{enumerate}
\def\labelenumi{\arabic{enumi}.}
\tightlist
\item
  a complete flatness check for finite-group connections on square
  lattices, with tests (Section 2);
\item
  an exact replay of the five retained C017 lattice cases and of every
  retained model solution record for them, with per-arm counts (Section
  3);
\item
  a table of the retained formal-mathematics pilot runs for challenges
  C045 to C050, with what the paired checkers accepted, what the prompts
  disclosed and the limits the harness set (Section 4).
\end{enumerate}

No new model runs were made for this revision. The earlier internal
revision numbered 2.0.1 added a replay, with no model call, of answers
that need no mathematics against the C045 to C050 checkers as they were
when the runs were recorded (Section 4.2).

The initial internal draft of this paper (2026-09-28) also described a
five-stage prompt filter whose stages were named after immune
mechanisms, and it reported a 50 of 50 benchmark, a 40 per cent
reduction in premature rollbacks, an 83.5 per cent reduction in prompt
tokens and a 50 per cent accuracy gain. No run, runner or receipt for
any of these is retained in the repository, so they are withdrawn; the
change log lists every removed statement. In the initial internal
draft's package the stages were simple functions. ``MHC peptide
decomposition'' split an edge list into radius-one neighbourhoods.
``Thymic selection'' tested membership in a supplied list. ``Kinetic
proofreading'' compared two supplied fractions with two thresholds. The
``Treg scratchpad'' check tested that four labelled lines appear in
order. The ``allosteric quorum'' counted how many of five listed Unicode
symbols occur in the prompt and compared the count with three. No effect
of any of them on model answers was measured.

\hypertarget{a-complete-flatness-check}{%
\subsection{2. A complete flatness
check}\label{a-complete-flatness-check}}

\hypertarget{definitions}{%
\subsubsection{2.1 Definitions}\label{definitions}}

A finite group \(G\) is given by named elements, a full multiplication
table and an inverse map. The package checks the identity law and the
inverse law for every element and associativity for every triple before
it uses a table, and it refuses a table that fails. It builds \(S_3\)
from the permutations of \(\{0,1,2\}\) and the quaternion group \(Q_8\)
from the unit quaternions \(\pm 1, \pm i, \pm j, \pm k\) with integer
coordinates.

An \(L \times M\) lattice has vertices \((r,c)\) with \(0 \le r < L\)
and \(0 \le c < M\). Its positive edges are \((r,c) \to (r,c+1)\) and
\((r,c) \to (r+1,c)\); on a torus the indices are taken modulo \(L\) and
\(M\). An open lattice requires \(L,M \ge 2\) in this implementation. A
torus requires \(L,M \ge 3\): at size two, distinct parallel edges would
have the same endpoint names and need a different representation. A
connection assigns an element \(U_{uv} \in G\) to each positive edge,
and a reversed edge carries the inverse, \(U_{vu} = U_{uv}^{-1}\). The
product along a path \(v_0 \to v_1 \to \cdots \to v_k\) is
\[\operatorname{hol}(v_0 \to \cdots \to v_k) = U_{v_0 v_1} U_{v_1 v_2} \cdots U_{v_{k-1} v_k}.\]
The plaquette \(P_{r,c}\) is listed as \(v_0 = (r,c)\),
\(v_1 = (r+1,c)\), \(v_2 = (r+1,c+1)\), \(v_3 = (r,c+1)\), and its
product is
\[\operatorname{hol}(P_{r,c}) = U_{v_0 v_1} U_{v_1 v_2} U_{v_2 v_3} U_{v_3 v_0}.\]
The last two factors traverse positive edges backwards, so they are the
inverses of the stored links. A connection is flat when every plaquette
product is the identity \(e\).

A gauge \(g\) assigns an element to each vertex and acts by
\(U'_{uv} = g_u U_{uv} g_v^{-1}\). The product around a closed path
based at \(v_0\) becomes \(g_{v_0} \operatorname{hol} g_{v_0}^{-1}\).
Flatness is therefore unchanged by a gauge transformation. On a torus,
let \(A\) be the product around
\((0,0) \to (0,1) \to \cdots \to (0,M-1) \to (0,0)\) and \(B\) the
product around \((0,0) \to (1,0) \to \cdots \to (L-1,0) \to (0,0)\). A
gauge transformation conjugates \(A\) and \(B\) by the same element.

We use four standard facts {[}1, 2, 3{]}. Reversing a closed path
inverts its product. Starting the same closed path at \(v_1\) instead of
\(v_0\) conjugates its product by \(U_{v_0 v_1}^{-1}\). If a connection
is flat, the product along a closed path depends only on the path's
homotopy class, so on an open lattice, which is simply connected, a
gauge built along a spanning tree makes every link the identity. On the
torus, whose fundamental group \(\mathbb{Z}^2\) is abelian, a flat
connection has commuting \(A\) and \(B\), and it is gauge-trivial
exactly when \(A = B = e\).

\hypertarget{the-check}{%
\subsubsection{2.2 The check}\label{the-check}}

The checker in \texttt{code/\allowbreak{}holonomy.py} returns one of
three verdicts.

\begin{enumerate}
\def\labelenumi{\arabic{enumi}.}
\tightlist
\item
  It reads the group table and refuses it unless it is a group. It reads
  the links and refuses them unless every edge is listed in both
  directions with \(U_{vu} = U_{uv}^{-1}\) and no other edge is listed.
  Every value must be a declared group element; duplicate parsed edges
  are refused even if alternate strings such as
  \texttt{0,0-\allowbreak{}>0,1} and \texttt{00,0-\allowbreak{}>0,1}
  carry equal values.
\item
  It computes every plaquette product in the order of Section 2.1. If
  any is not \(e\), it returns \texttt{NOT\_\allowbreak{}FLAT} with the
  first such plaquette in row-major order and its product.
\item
  On an open lattice it returns \texttt{FLAT\_\allowbreak{}TRIVIAL} with
  the gauge \(g_{(0,0)} = e\), \(g_v = g_u U_{uv}\) along a
  breadth-first spanning tree, after verifying that this gauge makes
  every link the identity.
\item
  On a torus it computes \(A\) and \(B\) and stops with an error if they
  do not commute, which cannot happen for a flat connection. If
  \(A = B = e\) it returns \texttt{FLAT\_\allowbreak{}TRIVIAL} with the
  verified tree gauge; otherwise it returns
  \texttt{FLAT\_\allowbreak{}NONTRIVIAL} with \(A\) and \(B\).
\end{enumerate}

\texttt{FLAT\_\allowbreak{}TRIVIAL} is the verdict that challenge C017
calls \texttt{PURE\_\allowbreak{}GAUGE}.

A separate function checks a stated verdict and its certificate. For
\texttt{NOT\_\allowbreak{}FLAT} the four listed vertices must bound a
plaquette, in either orientation and from any starting vertex, and the
stated product must equal the product in the listed order and differ
from \(e\). For \texttt{FLAT\_\allowbreak{}NONTRIVIAL} every plaquette
must be flat and at least one stated closed loop must have the stated,
non-identity product. In a flat connection every contractible loop has
product \(e\), so such a loop is non-contractible. For
\texttt{FLAT\_\allowbreak{}TRIVIAL} the stated gauge must make every
link the identity. Any other answer, including an answer that could not
be parsed, is rejected with a reason.

A gauge certificate must cover each vertex once after parsing; an alias
may not silently replace a conflicting value. These checks concern
parsed mappings; raw JSON duplicate-key handling is the responsibility
of the caller's JSON decoder.

\textbf{Correctness and cost.} A nonidentity plaquette directly
witnesses non-flatness. If all plaquettes are identity, deleting an
immediate backtrack or replacing a path segment by the other sides of a
plaquette preserves its product. On the open rectangle these moves
contract every closed path, so root-to-vertex products are
path-independent and the tree gauge trivializes all edges. On the torus,
closed paths reduce to the two generators with the commutator relation:
their products commute, and all closed products are identity exactly
when both generators are. This proves the three-way decision for the
declared rectangular cell complexes using the standard facts of Section
2.1; it is not a formal proof of the Python implementation. Table
validation costs \(O(|G|^3)\) operations and \(O(|G|^2)\) table storage.
Once the group is validated, the lattice decision costs \(O(LM)\) table
operations and storage. A certificate check costs \(O(LM+k)\), where
\(k\) is the total number of submitted loop edges; malformed inputs
receive a typed refusal rather than a flatness verdict.

\hypertarget{remark-skipping-plaquettes-with-four-identity-links}{%
\subsubsection{2.3 Remark: skipping plaquettes with four identity
links}\label{remark-skipping-plaquettes-with-four-identity-links}}

If the four links of a plaquette are the identity, its product is the
identity. This uses only the identity law and \(e^{-1} = e\); the
package records the one-line proof in Lean
(\texttt{identity\_\allowbreak{}links\_\allowbreak{}give\_\allowbreak{}identity\_\allowbreak{}product}
in \texttt{lean4/\allowbreak{}LatticeDecimation.lean}). The initial
internal draft called this observation a theorem. It is a direct
consequence of the group axioms.

The checker may skip such plaquettes. It then performs fewer
multiplications and reaches the same result, which the tests confirm on
every connection they build. The saving depends on the gauge. If
\(U = e\) and \(g\) is a gauge, then \(U'_{uv} = g_u g_v^{-1}\) is the
identity only when \(g_u = g_v\). A gauge that differs on neighbouring
vertices therefore leaves no plaquette with four identity links, while
the connection stays flat and gauge-trivial. One retained C017 case,
\texttt{base\_\allowbreak{}02}, has this form: 28 of its 31 links are
not the identity, none of its 12 plaquettes can be skipped, and it is
gauge-trivial. Skipping is an optimization. It cannot replace the
products of the remaining plaquettes, and it cannot show flatness on its
own.

\hypertarget{tests}{%
\subsubsection{2.4 Tests}\label{tests}}

The tests in
\texttt{code/\allowbreak{}tests/\allowbreak{}test\_\allowbreak{}holonomy\_\allowbreak{}checker.py}
check the following on \(S_3\) and \(Q_8\):

\begin{itemize}
\tightlist
\item
  the trivial connection is \texttt{FLAT\_\allowbreak{}TRIVIAL} on open
  and periodic lattices;
\item
  a connection that is a gauge transformation of the trivial one, with
  no plaquette of four identity links, is
  \texttt{FLAT\_\allowbreak{}TRIVIAL};
\item
  a torus connection with a constant non-identity product around one
  cycle is \texttt{FLAT\_\allowbreak{}NONTRIVIAL};
\item
  a single non-identity link away from the boundary gives
  \texttt{NOT\_\allowbreak{}FLAT} with exactly two non-flat plaquettes;
\item
  for every pair \((a,b) \in G \times G\), the \(3 \times 4\) torus
  connection whose links from the last column to the first carry \(a\),
  whose links from the last row to the first carry \(b\), and whose
  other links are \(e\), is flat exactly when \(ab = ba\); when it is
  flat its cycle products are \((a, b)\), or their conjugates by one
  element after a gauge transformation, and when it is not flat its only
  non-flat plaquette has product \(b a b^{-1} a^{-1}\);
\item
  on seeded random connections, reversing a plaquette inverts its
  product, changing the starting vertex conjugates it, and some products
  differ from their inverses;
\item
  skipping plaquettes with four identity links never changes the result;
\item
  malformed group tables, lattices and link lists are refused, and the
  certificate check accepts a correct plaquette in either orientation
  and rejects a wrong product, a flat plaquette, a four-cycle that is
  not a plaquette and a contractible loop.
\end{itemize}

The seam family is small, but it is exhaustive over \(G \times G\), and
in it a flat connection with non-commuting cycle products never
occurred, as the theory requires. The family contains 100 connections:
18 commuting S3 pairs and 40 commuting Q8 pairs. It tests 58 flat and 42
non-flat cases, not all connections on a lattice.

\hypertarget{the-c017-pilot}{%
\subsection{3. The C017 pilot}\label{the-c017-pilot}}

\hypertarget{what-was-run}{%
\subsubsection{3.1 What was run}\label{what-was-run}}

Challenge C017 asks for \texttt{NOT\_\allowbreak{}FLAT},
\texttt{PURE\_\allowbreak{}GAUGE} or
\texttt{FLAT\_\allowbreak{}NONTRIVIAL} verdicts with certificates for
connections with values in \(S_3\), \(Q_8\) and
\(\mathrm{SL}(2, \mathbb{F}_3)\) on square lattices. Its specification
calls for open-weight models of at most 70B parameters and 36 hidden
base cases. The retained runs are smaller pilots. They used five cases
from the challenge's judge packet (Table 1), one reply per case, and one
model outside that specification. The five case files are copied into
\texttt{code/\allowbreak{}data/\allowbreak{}c017/\allowbreak{}} with
their SHA-256 digests and source paths.

Each arm sent one prompt per case and parsed a JSON object from the
reply. Two checkers in the harness scored the certificate, and a case
passed only when both accepted it.

\begin{itemize}
\tightlist
\item
  \textbf{Arm 1.} The lattice, the group table and inverses, all
  directed links, the four verdict definitions and the JSON format.
  Model
  \texttt{meta-\allowbreak{}llama/\allowbreak{}llama-\allowbreak{}3.3-\allowbreak{}70b-\allowbreak{}instruct};
  temperature 0.1; at most 2,000 completion tokens.
\item
  \textbf{Arm 2.} Arm 1 plus four written constraints: multiply in
  traversal order; test the product, not a sign or determinant; on a
  torus a flat connection need not be gauge-trivial; refuse inverse
  mismatches. Same model and settings.
\item
  \textbf{Arm 5.} The Arm 2 constraints plus the full list of plaquettes
  in the vertex order that the checker accepts, an instruction to
  multiply step by step by table lookup, the two cycle paths on a torus
  and a decision procedure. Run with
  \texttt{meta-\allowbreak{}llama/\allowbreak{}llama-\allowbreak{}3.3-\allowbreak{}70b-\allowbreak{}instruct}
  (temperature 0.05, at most 3,500 completion tokens) and with
  \texttt{stealth/\allowbreak{}space-\allowbreak{}bunny-\allowbreak{}alpha}
  (temperature 0.05; the runner's default limit is 6,000 completion
  tokens, can be changed on the command line, and is not recorded in the
  report).
\item
  \textbf{Arm 6.} Arm 5, except that the harness listed only the
  plaquettes that contain a non-identity link and stated that the others
  were flat. Run with
  \texttt{stealth/\allowbreak{}space-\allowbreak{}bunny-\allowbreak{}alpha}
  (temperature 0.05; runner default limit 4,000 completion tokens).
\end{itemize}

No C017 runs are retained for Arms 3 and 4. The model
\texttt{stealth/\allowbreak{}space-\allowbreak{}bunny-\allowbreak{}alpha}
is a new free ``stealth'' model offered on OpenRouter as of 2026-09-28;
the provider does not disclose its identity, size or training. Results
with it are exploratory. They fall outside the challenge protocol, SG-CP
v1, whose §2 admits only open-weight models of at most 70B parameters
(\texttt{omega\_\allowbreak{}engine\_\allowbreak{}projects/\allowbreak{}mssn\_\allowbreak{}sn\_\allowbreak{}symgam/\allowbreak{}challenges/\allowbreak{}01\_\allowbreak{}challenge\_\allowbreak{}protocol.md}),
and they may not be reproducible once the stealth period ends. The
reports record costs computed by the harness from configured per-token
prices, not provider bills, so we do not report costs.

\Needspace{18\baselineskip}

Table 1. The five retained C017 cases and the host verdicts computed by
the package. ``Skippable'' counts plaquettes whose four links are the
identity. The case files name the elements of \(S_3\) as \texttt{e},
\texttt{r}, \texttt{r2}, \texttt{s}, \texttt{sr}, \texttt{sr2}, where
\texttt{r2} is \(r^2 = r^{-1}\), and the elements of \(Q_8\) as
\texttt{e}, \texttt{-\allowbreak{}e}, \texttt{i},
\texttt{-\allowbreak{}i}, \texttt{j}, \texttt{-\allowbreak{}j},
\texttt{k}, \texttt{-\allowbreak{}k}. The case
\texttt{pois\_\allowbreak{}01} comes from the challenge's poisoned set:
its non-flat product is a 3-cycle, which has sign \(+1\), so a test of
the sign alone would call the connection flat.

\begin{longtable}[]{@{}
  >{\raggedright\arraybackslash}p{(\columnwidth - 8\tabcolsep) * \real{0.1200}}
  >{\raggedright\arraybackslash}p{(\columnwidth - 8\tabcolsep) * \real{0.2400}}
  >{\raggedright\arraybackslash}p{(\columnwidth - 8\tabcolsep) * \real{0.1800}}
  >{\raggedright\arraybackslash}p{(\columnwidth - 8\tabcolsep) * \real{0.2000}}
  >{\raggedright\arraybackslash}p{(\columnwidth - 8\tabcolsep) * \real{0.2600}}@{}}
\toprule\noalign{}
\begin{minipage}[b]{\linewidth}\raggedright
Case
\end{minipage} & \begin{minipage}[b]{\linewidth}\raggedright
Group and lattice
\end{minipage} & \begin{minipage}[b]{\linewidth}\raggedright
Non-identity links
\end{minipage} & \begin{minipage}[b]{\linewidth}\raggedright
Plaquettes (skippable)
\end{minipage} & \begin{minipage}[b]{\linewidth}\raggedright
Host verdict
\end{minipage} \\
\midrule\noalign{}
\endhead
\bottomrule\noalign{}
\endlastfoot
\texttt{pois\_\allowbreak{}01} & \(S_3\), \(3 \times 3\) open & 1 of 12
& 4 (3) & \texttt{NOT\_\allowbreak{}FLAT}, \(P_{0,0}\) has product
\texttt{r2} \\
\texttt{base\_\allowbreak{}01} & \(S_3\), \(3 \times 4\) open & 1 of 17
& 6 (5) & \texttt{NOT\_\allowbreak{}FLAT}, \(P_{0,1}\) has product
\texttt{s} \\
\texttt{base\_\allowbreak{}02} & \(S_3\), \(4 \times 5\) open & 28 of 31
& 12 (0) & \texttt{FLAT\_\allowbreak{}TRIVIAL} \\
\texttt{base\_\allowbreak{}03} & \(S_3\), \(5 \times 6\) torus & 11 of
60 & 30 (20) & \texttt{FLAT\_\allowbreak{}NONTRIVIAL}, \(A\) =
\texttt{r}, \(B\) = \texttt{r2} \\
\texttt{base\_\allowbreak{}13} & \(Q_8\), \(4 \times 5\) open & 1 of 31
& 12 (11) & \texttt{NOT\_\allowbreak{}FLAT}, \(P_{0,0}\) has product
\texttt{-\allowbreak{}i} \\
\end{longtable}

The host verdicts agree with the verdicts that the harness expected for
all five cases. The \(S_3\) table in the case files is isomorphic to the
permutation group (the tests find all six isomorphisms), and the \(Q_8\)
table equals the package's quaternion table.

\hypertarget{outcomes}{%
\subsubsection{3.2 Outcomes}\label{outcomes}}

Table 2 gives, for each arm and case, what the model stated and whether
the harness accepted the certificate. ``70B'' denotes
\texttt{meta-\allowbreak{}llama/\allowbreak{}llama-\allowbreak{}3.3-\allowbreak{}70b-\allowbreak{}instruct}
and ``stealth'' denotes
\texttt{stealth/\allowbreak{}space-\allowbreak{}bunny-\allowbreak{}alpha}.
``Reversed order'' means that the answer named the correct non-flat
plaquette with its vertices in the opposite orientation.

\Needspace{20\baselineskip}

Table 2. Harness outcomes for the five C017 cases (one reply per cell).

\begin{longtable}[]{@{}
  >{\raggedright\arraybackslash}p{(\columnwidth - 10\tabcolsep) * \real{0.1667}}
  >{\raggedright\arraybackslash}p{(\columnwidth - 10\tabcolsep) * \real{0.1667}}
  >{\raggedright\arraybackslash}p{(\columnwidth - 10\tabcolsep) * \real{0.1667}}
  >{\raggedright\arraybackslash}p{(\columnwidth - 10\tabcolsep) * \real{0.1667}}
  >{\raggedright\arraybackslash}p{(\columnwidth - 10\tabcolsep) * \real{0.1667}}
  >{\raggedright\arraybackslash}p{(\columnwidth - 10\tabcolsep) * \real{0.1667}}@{}}
\toprule\noalign{}
\begin{minipage}[b]{\linewidth}\raggedright
Case
\end{minipage} & \begin{minipage}[b]{\linewidth}\raggedright
Arm 1, 70B
\end{minipage} & \begin{minipage}[b]{\linewidth}\raggedright
Arm 2, 70B
\end{minipage} & \begin{minipage}[b]{\linewidth}\raggedright
Arm 5, 70B
\end{minipage} & \begin{minipage}[b]{\linewidth}\raggedright
Arm 5, stealth
\end{minipage} & \begin{minipage}[b]{\linewidth}\raggedright
Arm 6, stealth
\end{minipage} \\
\midrule\noalign{}
\endhead
\bottomrule\noalign{}
\endlastfoot
\texttt{pois\_\allowbreak{}01} & reversed order & reversed order &
accepted & accepted & accepted \\
\texttt{base\_\allowbreak{}01} & reversed order & stated
\texttt{PURE\_\allowbreak{}GAUGE} & accepted & accepted & accepted \\
\texttt{base\_\allowbreak{}02} & no parsable JSON & stated
\texttt{NOT\_\allowbreak{}FLAT} & flat plaquette stated as \texttt{r2} &
flat plaquette stated as \texttt{r} & flat plaquette stated as
\texttt{r} \\
\texttt{base\_\allowbreak{}03} & loop product misstated & loop product
misstated & no parsable JSON & reply cut off & accepted \\
\texttt{base\_\allowbreak{}13} & reversed order & no parsable JSON &
accepted & accepted & accepted \\
Accepted & 0 of 5 & 0 of 5 & 3 of 5 & 3 of 5 & 4 of 5 \\
Verdict label equal to host & 4 of 5 & 2 of 5 & 3 of 5 & 3 of 5 & 4 of
5 \\
\end{longtable}

For the 15 solution records retained in Arms 5 and 6, the package agrees
with the historical harness: 10 accepted and 5 rejected. These records
contain 13 parsed certificate objects and two
\texttt{PARSE\_\allowbreak{}ERROR} wrappers for replies whose JSON could
not be parsed; rejecting a wrapper is not verification of an absent
mathematical certificate. For Arms 1 and 2 the reports retain only the
stated verdict and the checkers' reasons, not the certificates.

\textbf{Vertex order.} The Arm 1 and Arm 2 prompts defined a plaquette
product as \(U_{01} U_{12} U_{23} U_{30}\) but did not say in which
order the checker lists a plaquette's vertices. The checker accepted
only \((r,c), (r+1,c), (r+1,c+1), (r,c+1)\). Three Arm 1 answers named
the correct non-flat plaquette in the order
\((r,c), (r,c+1), (r+1,c+1), (r+1,c)\) and were rejected on that ground
before any product was compared. The reports do not retain the stated
products, so we cannot say whether they were right for the reversed
order. Arm 1 therefore scored 0 of 5, although its verdict label matched
the host in 4 of 5 cases. Two Arm 2 answers were rejected on the same
ground; one of them also had the wrong verdict. The Arm 5 prompt listed
the accepted order, and the three open non-flat cases were accepted for
both models. Arm 5 also added the step-by-step instruction and the cycle
paths, so these runs do not separate the effect of stating the
convention from the other additions. The certificate check of Section 2
accepts either orientation.

\textbf{The gauge-trivial case.} \texttt{base\_\allowbreak{}02} is flat
and gauge-trivial, and 28 of its 31 links are not the identity. All
three Arm 5 and Arm 6 answers for it stated
\texttt{NOT\_\allowbreak{}FLAT} with a plaquette whose product is \(e\),
stating the product as \texttt{r2} or \texttt{r}; each was rejected
because the stated product disagrees with the table. The initial
internal draft also stated that an instruction to build a gauge along a
spanning tree biased models towards \texttt{PURE\_\allowbreak{}GAUGE}.
No retained run supports this, and the statement is withdrawn. The only
retained \texttt{PURE\_\allowbreak{}GAUGE} answer is the Arm 2 answer
for \texttt{base\_\allowbreak{}01}, a non-flat case, whose stated gauge
fails on the edge \((0,1) \to (0,2)\).

\textbf{The torus case.} In \texttt{base\_\allowbreak{}03}, 20 of the 30
plaquettes have four identity links, and each of the other 10 contains
at least one of the 11 non-identity links. With Arm 5, the reply of
\texttt{stealth/\allowbreak{}space-\allowbreak{}bunny-\allowbreak{}alpha}
stopped at 8,000 completion tokens with its JSON unfinished; the
retained start of the reply already carried the label
\texttt{FLAT\_\allowbreak{}NONTRIVIAL}. With Arm 6, where the harness
listed only the 10 plaquettes with a non-identity link, the same model
returned all 10 products (each \(e\)), \(A\) = \texttt{r}, \(B\) =
\texttt{r2} and the label \texttt{FLAT\_\allowbreak{}NONTRIVIAL} in
2,810 completion tokens, and both checkers accepted the answer. This is
one case with one reply per arm. It shows that the shorter prompt fitted
within the limit on this case; it does not measure a general saving. The
Arm 6 listing relies on identity links, which a gauge transformation can
remove (Section 2.3). The initial internal draft stated that 24
plaquettes were skipped and that the model evaluated only the two
cycles. Recomputing the Arm 6 plaquette list from the case file gives 20
skipped plaquettes, and the retained reply evaluates the other 10
plaquettes and the two cycles.

\textbf{The harness checkers.} The two checkers in the C017 harness are
not independent implementations. Both import the plaquette list and the
path product from the challenge's reference solver. The check in this
package shares no code with them, and on the retained solution records
it reaches the same decisions. Agreement on these cases does not certify
independence from shared mathematical assumptions or establish the
soundness of either implementation beyond the checked scope.

\hypertarget{data-sources-for-section-3}{%
\subsubsection{3.3 Data sources for Section
3}\label{data-sources-for-section-3}}

All paths are relative to the repository root.

\begin{itemize}
\tightlist
\item
  Challenge specification:
  \texttt{omega\_\allowbreak{}engine\_\allowbreak{}projects/\allowbreak{}mssn\_\allowbreak{}sn\_\allowbreak{}symgam/\allowbreak{}challenges/\allowbreak{}b\_\allowbreak{}cohomology/\allowbreak{}c017\_\allowbreak{}nonabelian\_\allowbreak{}lattice\_\allowbreak{}holonomy.md}.
\item
  Arms 1 and 2:
  \texttt{.../\allowbreak{}c017\_\allowbreak{}nonabelian\_\allowbreak{}lattice\_\allowbreak{}holonomy/\allowbreak{}results/\allowbreak{}immune\_\allowbreak{}transfer\_\allowbreak{}report.json},
  added in commit
  \texttt{\seqsplit{7e639a0be573f6affc4d84a2aacb5573da92bc79}}; runner
  \texttt{eval\_\allowbreak{}epistemic\_\allowbreak{}immune\_\allowbreak{}transfer.py}
  at that commit.
\item
  Arm 5 with
  \texttt{meta-\allowbreak{}llama/\allowbreak{}llama-\allowbreak{}3.3-\allowbreak{}70b-\allowbreak{}instruct}:
  \texttt{results/\allowbreak{}arm5\_\allowbreak{}hybrid\_\allowbreak{}report.json}
  as committed in
  \texttt{\seqsplit{9e45ee501e6a48086aaabd2381632eef6f7ed7a2}}. The file
  at that path was later overwritten by a one-case run, so this report
  is retained only in the repository history.
\item
  Arm 5 with
  \texttt{stealth/\allowbreak{}space-\allowbreak{}bunny-\allowbreak{}alpha}:
  \texttt{results/\allowbreak{}arm5\_\allowbreak{}stealth\_\allowbreak{}space\_\allowbreak{}bunny\_\allowbreak{}alpha\_\allowbreak{}report.json},
  commit \texttt{\seqsplit{5327cfcb7ee3622e79f7d0cd884dfa37ef4eef56}}.
\item
  Arm 6:
  \texttt{results/\allowbreak{}arm6\_\allowbreak{}stealth\_\allowbreak{}space\_\allowbreak{}bunny\_\allowbreak{}alpha\_\allowbreak{}report.json},
  commit \texttt{\seqsplit{eb5fc971ae44320ff8183f09343d547f830d9d56}}.
\end{itemize}

The historical report digests are recorded with their original receipts.
The byte copies in
\texttt{docs/\allowbreak{}historical\_\allowbreak{}evidence/\allowbreak{}}
have pinned source commits, Git blob SHAs and SHA-256 digests in their
provenance file. The package file
\texttt{code/\allowbreak{}data/\allowbreak{}c017/\allowbreak{}model\_\allowbreak{}claims.json}
holds the verdicts, certificates, checker reasons and completion lengths
extracted from these reports.
\texttt{code/\allowbreak{}recount\_\allowbreak{}pilots.py} freshly
compares all 25 extraction entries with those source reports. That
comparison does not recover the ten missing certificate payloads or
generate new responses.

\hypertarget{six-formal-mathematics-pilots}{%
\subsection{4. Six formal-mathematics
pilots}\label{six-formal-mathematics-pilots}}

\hypertarget{what-was-run-1}{%
\subsubsection{4.1 What was run}\label{what-was-run-1}}

For each of challenges C045 to C050, a runner script defines a ten-case
panel, separate from the challenge's 64 hidden cases (for C049 panel B
the script is not retained; see below), and sends each case once to
three models through OpenRouter at temperature 0:
\texttt{stealth/\allowbreak{}space-\allowbreak{}bunny-\allowbreak{}alpha}
(the undisclosed stealth model described in Section 3.1),
\texttt{meta-\allowbreak{}llama/\allowbreak{}llama-\allowbreak{}3.3-\allowbreak{}70b-\allowbreak{}instruct}
and
\texttt{meta-\allowbreak{}llama/\allowbreak{}llama-\allowbreak{}3.1-\allowbreak{}8b-\allowbreak{}instruct}.
The reports do not record which provider served each model. A case
passed only when both of the challenge's checkers, an exact checker and
a second checker, accepted the answer.

The harness set the completion limits. In C045 and C046 the limit was
4,096 tokens for
\texttt{stealth/\allowbreak{}space-\allowbreak{}bunny-\allowbreak{}alpha}
and 2,048 for the other two models. In C047 and C049 panel A it was
4,096 tokens and in C048 and C050 2,048 tokens for all three. The limit
for C049 panel B is not recorded.

\Needspace{25\baselineskip}

Table 3. Cases passed, out of 10 per cell, one reply per case. SBA, 70B
and 8B refer to the three model identifiers above, in that order.

\begin{longtable}[]{@{}
  >{\raggedright\arraybackslash}p{(\columnwidth - 8\tabcolsep) * \real{0.1300}}
  >{\raggedright\arraybackslash}p{(\columnwidth - 8\tabcolsep) * \real{0.4200}}
  >{\raggedright\arraybackslash}p{(\columnwidth - 8\tabcolsep) * \real{0.1500}}
  >{\raggedright\arraybackslash}p{(\columnwidth - 8\tabcolsep) * \real{0.1500}}
  >{\raggedright\arraybackslash}p{(\columnwidth - 8\tabcolsep) * \real{0.1500}}@{}}
\toprule\noalign{}
\begin{minipage}[b]{\linewidth}\raggedright
Challenge
\end{minipage} & \begin{minipage}[b]{\linewidth}\raggedright
Task
\end{minipage} & \begin{minipage}[b]{\linewidth}\raggedright
SBA
\end{minipage} & \begin{minipage}[b]{\linewidth}\raggedright
70B
\end{minipage} & \begin{minipage}[b]{\linewidth}\raggedright
8B
\end{minipage} \\
\midrule\noalign{}
\endhead
\bottomrule\noalign{}
\endlastfoot
C045 & knot invariants from planar diagrams & 2 & 8 & 3 \\
C046 & Knuth--Bendix completion and normal forms {[}4{]} & 4 & 7 & 8 \\
C047 & linear recurrences or certified absence & 4 & 2 & 0 \\
C048 & rational point or local obstruction for a ternary quadratic form
& 10 & 8 & 5 \\
C049, panel A & limits, colimits and adjunctions in finite categories &
5 & 6 & 4 \\
C049, panel B & the same topics, a different set of ten cases & 9 & 6 &
6 \\
C050, run 1 & equational rewriting proof or two-element countermodel & 7
& 4 & 1 \\
C050, run 2 & the same ten cases, run again & 6 & 3 & 1 \\
\end{longtable}

The offline recount uses the per-case recorded flags. One C045 8B entry
records \texttt{JSON\_\allowbreak{}ERROR} without a pass flag; it is
counted as an unaccepted response. This recount does not re-execute the
historical certificate checkers.

C049 has two retained runs on different ten-case panels with the same
three models, committed together on 2026-09-28. Panel A was produced by
the runner script retained in the repository; the script that produced
panel B is not retained. For
\texttt{stealth/\allowbreak{}space-\allowbreak{}bunny-\allowbreak{}alpha}
the two panels give 5 and 9 passes. C050 was run twice with the same
committed runner and the same ten cases; the recorded start times are 98
seconds apart. The second report overwrote the first, which is retained
in the repository history. For
\texttt{stealth/\allowbreak{}space-\allowbreak{}bunny-\allowbreak{}alpha},
five of the ten cases changed outcome between the two runs, two of them
because one run received an API error instead of a reply. The initial
internal draft reported only C049 panel B and C050 run 2.

\textbf{What the prompts disclosed.} Each runner sent a case as the line
\texttt{Case ID: <id>} followed by the case data (lines 219, 270, 221,
163, 264 and 236 of the C045 to C050 runners, unchanged from the commits
that recorded the runs to the code baseline). The challenge authors
wrote the identifiers, and many name the expected result.
\texttt{c045\_\allowbreak{}m02\_\allowbreak{}trefoil\_\allowbreak{}identify}
asks for the identification of a knot that its identifier names.
\texttt{c045\_\allowbreak{}m09\_\allowbreak{}refusal\_\allowbreak{}corrupt\_\allowbreak{}edge}
and
\texttt{c046\_\allowbreak{}m08\_\allowbreak{}refusal\_\allowbreak{}commutativity}
name a refusal and its reason.
\texttt{c046\_\allowbreak{}m05\_\allowbreak{}check\_\allowbreak{}non\_\allowbreak{}confluent\_\allowbreak{}group}
names the confluence verdict.
\texttt{c047\_\allowbreak{}m07\_\allowbreak{}prand\_\allowbreak{}decoy}
and
\texttt{c047\_\allowbreak{}m08\_\allowbreak{}order9\_\allowbreak{}decoy}
mark the cases with no recurrence within the bounds.
\texttt{c048\_\allowbreak{}m01\_\allowbreak{}solv\_\allowbreak{}11\_\allowbreak{}m2}
and
\texttt{c048\_\allowbreak{}m02\_\allowbreak{}solv\_\allowbreak{}pyth}
mark solvable forms, and
\texttt{c048\_\allowbreak{}m05\_\allowbreak{}obstr\_\allowbreak{}inf\_\allowbreak{}111}
names both the obstruction and its place.
\texttt{c049\_\allowbreak{}m05\_\allowbreak{}no\_\allowbreak{}product\_\allowbreak{}disc2}
names the verdict that no product exists. The C045, C046, C049 and C050
runners also sent the whole case record, including the \texttt{note} of
the poisoned cases C049 m10 and C050 m09, whose identifiers call the
note bait or fake; the C047 and C048 runners sent only the terms or
coefficients and the bounds. Table 4 counts the identifiers of the six
panels whose runner is retained. Which identifiers name the result is
one reader's classification, the author's; it is published identifier by
identifier with the replay listed in Section 4.4, so that another reader
can check or dispute each case.

\Needspace{16\baselineskip}

Table 4. What the case identifiers name, in the author's classification,
and the most passes that an answer needing no mathematics obtains under
the same rule (both checkers accept), out of 10 per panel. The answers
tried were an empty object, a bare abstention, a refusal or out-of-scope
answer with each reason the prompt lists, and a constant verdict where
the output format allows one (Section 4.2).

\begin{longtable}[]{@{}
  >{\raggedright\arraybackslash}p{(\columnwidth - 8\tabcolsep) * \real{0.1579}}
  >{\raggedleft\arraybackslash}p{(\columnwidth - 8\tabcolsep) * \real{0.2105}}
  >{\raggedleft\arraybackslash}p{(\columnwidth - 8\tabcolsep) * \real{0.2105}}
  >{\raggedleft\arraybackslash}p{(\columnwidth - 8\tabcolsep) * \real{0.2105}}
  >{\raggedleft\arraybackslash}p{(\columnwidth - 8\tabcolsep) * \real{0.2105}}@{}}
\toprule\noalign{}
\begin{minipage}[b]{\linewidth}\raggedright
Challenge
\end{minipage} & \begin{minipage}[b]{\linewidth}\raggedleft
Names the verdict or answer
\end{minipage} & \begin{minipage}[b]{\linewidth}\raggedleft
Names the object
\end{minipage} & \begin{minipage}[b]{\linewidth}\raggedleft
Neither
\end{minipage} & \begin{minipage}[b]{\linewidth}\raggedleft
Best baseline passes
\end{minipage} \\
\midrule\noalign{}
\endhead
\bottomrule\noalign{}
\endlastfoot
C045 & 7 & 3 & 0 & 10 \\
C046 & 5 & 3 & 2 & 10 \\
C047 & 4 & 6 & 0 & 0 \\
C048 & 10 & 0 & 0 & 3 \\
C049, panel A & 5 & 5 & 0 & 2 \\
C050 & 5 & 5 & 0 & 1 \\
\end{longtable}

In C045 and C046 the best such answer is a refusal of every case: both
checkers accept a refusal with any of the listed reasons on every case,
so it passes 10 of 10, more than any model in Table 3. In C048 a
constant local obstruction at the place 2 passes 3 cases, and in C049 a
constant ``does not exist'' passes 2. The identifiers of the C049 panel
B report have the same form (for example
\texttt{c049\_\allowbreak{}m07\_\allowbreak{}no\_\allowbreak{}product\_\allowbreak{}discrete}),
but its runner is not retained, so what its prompts contained cannot be
checked and its cases cannot be replayed.

The counts in Table 3 are therefore counts under leaked labels. These
counts describe responses under the disclosed prompts and checker rules;
they do not estimate or bound unaided performance. In C045 and C046 they
do not separate any model from an answer that refuses every case. The
models' own passes there include refusals only on the cases whose
identifiers name a refusal (C045 m09 and m10; C046 m08 to m10), and the
checkers accepted those refusals without comparing the stated reason
with the input.

We do not add counts across challenges and we make no comparison between
models. The tasks differ, the panels were written by the challenge
authors, each cell has n = 10 and one reply per case, and no statistical
test was run. The initial internal draft's totals over 80 cases and its
conclusion that absence proofs remain a capability of symbolic
algorithms alone are withdrawn; the evidence is 10 cases per cell on one
day.

\hypertarget{observations}{%
\subsubsection{4.2 Observations}\label{observations}}

\textbf{Harness limits.} In C045 all 8 failures of
\texttt{stealth/\allowbreak{}space-\allowbreak{}bunny-\allowbreak{}alpha}
returned no answer object, with recorded token counts (prompt and
completion together) of 4,751 to 4,845, consistent with the reply
reaching the 4,096-token completion limit. In C046, 5 of its 6 failures
were of the same kind, at 4,624 to 4,747 tokens. The reports do not
record prompt and completion counts separately. These cases most likely
show replies that did not fit the harness limit; they do not measure
what the model could compute.

\textbf{The second checkers.} For those empty answers in C045 and C046
the second checker reported acceptance and only the exact checker
rejected them. In C047 case m06, whose terms after the first are the
Fibonacci numbers,
\texttt{stealth/\allowbreak{}space-\allowbreak{}bunny-\allowbreak{}alpha}
returned a certificate that no recurrence exists within the bounds; the
second checker accepted it, and the exact checker rejected it because
its minor used rows below the smallest index that the exact checker
allows.

A replay of answers that need no mathematics, against the checkers as
they were when the runs were recorded, shows how far this goes. It
imports each challenge's two checkers from those commits, reads the
panel from the runner without running it, and calls no model. The second
checkers of C045 and C046 accept the empty object \texttt{\{\}}, a bare
abstention, and any refusal that carries a reason and a witness, on all
ten cases. The second checker of C047 accepts a recurrence with an empty
coefficient list, and an absence certificate with an empty list of minor
rows, on 9 of the 10 cases, because the determinant of an empty matrix
is 1; on the tenth, whose terms include a non-integer, it raises an
error. The second checker of C048 accepts any out-of-scope answer on all
ten cases, and that of C049 accepts a constant ``does not exist'' on 8
of 10. A pass required both checkers. For the cited empty or malformed
answers rejected by the exact checker, permissive acceptance by the
second checker did not turn a rejection into a pass. This does not prove
that the second checker adds no constraint on every possible answer.
Table 3 remains a historical record of the original paired rule. In C045
and C046 even that rule accepts the constant refusals of Table 4.

The repository documents later harness corrections in commits
\texttt{\seqsplit{e6fbc15cf011121597a82a1d199626063bedf7c0}} and
\texttt{\seqsplit{e2f7a3500413af960c6001832a1ab096a7a14af8}}. The
following is a historical summary of those changes, not a fresh checker
execution in this review. The C045 and C046 second checkers now
recompute the answer independently, so a refusal of every case,
\texttt{\{\}} and an answer without content fail them; the C047 second
checker no longer accepts an empty coefficient list or an empty list of
minor rows; the C048 second checker accepts an out-of-scope answer only
when a coefficient is not an integer; and the C043 to C050 runners
replace the descriptive identifiers with opaque tokens drawn for each
run. The C045 and C046 exact checkers still accept a refusal with any
reason, so the pair now depends on the second checker for that case.
None of this changes Table 3, which was produced with the checkers of
the run commits.

\textbf{Absence certificates.} C047 certifies that no recurrence exists
within the bounds by a nonsingular minor, whose determinant the exact
checker computes by fraction-free elimination {[}5{]}. On the three
panel cases that the runner marks as having no recurrence within the
bounds (m07, m08, m10), no model returned a certificate that both
checkers accepted.
\texttt{stealth/\allowbreak{}space-\allowbreak{}bunny-\allowbreak{}alpha}
returned no answer content for all three; the other two models stated
recurrences or malformed certificates.

\textbf{Other failures.} In each C050 run, one failure of
\texttt{stealth/\allowbreak{}space-\allowbreak{}bunny-\allowbreak{}alpha}
was an API error with no reply. In C048, the two failures of
\texttt{llama-\allowbreak{}3.3-\allowbreak{}70b-\allowbreak{}instruct}
were integer points that do not satisfy the form's equation.

\hypertarget{what-the-checkers-and-mutation-suites-show}{%
\subsubsection{4.3 What the checkers and mutation suites
show}\label{what-the-checkers-and-mutation-suites-show}}

Each challenge's judge packet retains 64 hidden cases and 48 mutated
certificates for testing the checkers, 512 and 384 in total over C043 to
C050. Scoring receipts are retained for C043 and C044 only. Each records
the reference solver passing all 64 hidden cases, the two checkers
agreeing on all of them, and all 48 mutations rejected. No scoring
receipt is retained for C045 to C050, so the initial internal draft's
statement that 512 cases were verified and 384 mutations rejected is
narrowed to these 128 cases and 96 mutations. The receipts concern the
reference solvers and checkers, not the models.

Challenges C043 and C044 were also run with the same three models. The
C043 runner wrote its report to a directory that the repository's ignore
rules exclude. The C044 runner wrote its report outside the repository,
to a directory it reaches six levels above its own file. Neither report
is under version control, and the counts appear only in narrative
findings documents, so we do not report them. They would not measure
identification in any case: the C044 runner sent each example file
whole, including its \texttt{note}, which names the group or the
expected verdict (for example ``Quaternion group Q8'', ``C4 against
Klein four group V4'' and ``S5 has order 120 which exceeds max allowed
order 64''). The C043 prompts were built from the case fields and
carried no note. An earlier internal review reported an external file
with the C044 report's name (SHA-256
\texttt{\seqsplit{5fecd63f08dc1b5b2094fe6b722c236f284567b52a9a140f62d4825e63e8648b}});
it was not under version control when the run was recorded. Its presence
and bytes were not independently checked in this review, and we do not
report its counts. In the corrected harness the C044 runner no longer
sends the \texttt{note} or reads a credentials file, and it writes its
report into the repository's run store.

\hypertarget{data-sources-for-section-4}{%
\subsubsection{4.4 Data sources for Section
4}\label{data-sources-for-section-4}}

All paths are relative to
\texttt{omega\_\allowbreak{}engine\_\allowbreak{}projects/\allowbreak{}mssn\_\allowbreak{}sn\_\allowbreak{}symgam/\allowbreak{}challenges/\allowbreak{}e\_\allowbreak{}formal\_\allowbreak{}math/\allowbreak{}}.

\begin{itemize}
\tightlist
\item
  C045 run report:
  \texttt{c045\_\allowbreak{}knot\_\allowbreak{}polynomials\_\allowbreak{}from\_\allowbreak{}pd/\allowbreak{}empirical\_\allowbreak{}results/\allowbreak{}c045\_\allowbreak{}space\_\allowbreak{}bunny\_\allowbreak{}matrix\_\allowbreak{}report.json},
  commit \texttt{\seqsplit{ea2b55d841bdbe283f2f6ac77abe8a1677cbf0f8}}.
\item
  C046, C047 and C048 run reports: the files with the same pattern in
  their challenge folders, commit
  \texttt{\seqsplit{46bd1dec156260f073b77caa07d8a453f8efad30}}.
\item
  C049 panel B:
  \texttt{c049\_\allowbreak{}finite\_\allowbreak{}category\_\allowbreak{}universal\_\allowbreak{}constructions/\allowbreak{}empirical\_\allowbreak{}results/\allowbreak{}c049\_\allowbreak{}space\_\allowbreak{}bunny\_\allowbreak{}matrix\_\allowbreak{}report.json};
  panel A: the three files \texttt{results\_\allowbreak{}*.json} in the
  same folder; both in commit
  \texttt{\seqsplit{46bd1dec156260f073b77caa07d8a453f8efad30}}.
\item
  C050 run 1:
  \texttt{c050\_\allowbreak{}equational\_\allowbreak{}theorem\_\allowbreak{}proving/\allowbreak{}empirical\_\allowbreak{}results/\allowbreak{}c050\_\allowbreak{}space\_\allowbreak{}bunny\_\allowbreak{}matrix\_\allowbreak{}report.json}
  as committed in
  \texttt{\seqsplit{46bd1dec156260f073b77caa07d8a453f8efad30}}; run 2:
  the same path as committed in
  \texttt{\seqsplit{5a256c6e8d088c612332362a8abc53138806ac7a}}.
\item
  Runner scripts, which set the panels, the prompts and the completion
  limits: \texttt{run\_\allowbreak{}c045\_\allowbreak{}matrix.py} to
  \texttt{run\_\allowbreak{}c050\_\allowbreak{}matrix.py} in the
  challenge folders.
\item
  The replay behind Table 4 and Section 4.2 (paths relative to the
  repository root): the collection's script
  \texttt{omega\_\allowbreak{}engine\_\allowbreak{}atlas/\allowbreak{}coda\_\allowbreak{}atlas/\allowbreak{}omega\_\allowbreak{}engine\_\allowbreak{}symbol\_\allowbreak{}game/\allowbreak{}tools/\allowbreak{}replay\_\allowbreak{}formal\_\allowbreak{}math\_\allowbreak{}baselines.py}
  and its record
  \texttt{omega\_\allowbreak{}engine\_\allowbreak{}atlas/\allowbreak{}coda\_\allowbreak{}atlas/\allowbreak{}omega\_\allowbreak{}engine\_\allowbreak{}symbol\_\allowbreak{}game/\allowbreak{}verification/\allowbreak{}formal\_\allowbreak{}math\_\allowbreak{}baselines.json}.
  It reads the runners and checkers from commits
  \texttt{\seqsplit{ea2b55d841bdbe283f2f6ac77abe8a1677cbf0f8}} (the C045
  runner) and
  \texttt{\seqsplit{46bd1dec156260f073b77caa07d8a453f8efad30}} (C046 to
  C050). The C045 checkers are read from commit
  \texttt{\seqsplit{c568c20062697b54ddb8729e2476d4cfdb5c4982}}, which
  differs from the run commit in that folder only by one import line
  that Python 3.12 needs to load the exact checker.
\item
  Scoring receipts:
  \texttt{c043\_\allowbreak{}ramsey\_\allowbreak{}vdw\_\allowbreak{}constructions/\allowbreak{}commit\_\allowbreak{}reveal/\allowbreak{}c043\_\allowbreak{}l41\_\allowbreak{}scoring\_\allowbreak{}receipt.json}
  (commit \texttt{\seqsplit{50197478726e5ecb791e519760cee9a79af2af27}})
  and
  \texttt{c044\_\allowbreak{}finite\_\allowbreak{}group\_\allowbreak{}identification/\allowbreak{}commit\_\allowbreak{}reveal/\allowbreak{}c044\_\allowbreak{}l41\_\allowbreak{}scoring\_\allowbreak{}receipt.json}
  (commit \texttt{\seqsplit{999b36c1afed0f905f634f83adc1c646da7b975e}}).
\end{itemize}

\hypertarget{limitations}{%
\subsection{5. Limitations}\label{limitations}}

The C017 pilot has five cases, one reply per case and arm, two models,
and arms that change several things at once. The formal-mathematics
pilots have ten cases per cell and one reply per case, and their prompts
named the expected result for 36 of the 60 cases whose prompts can be
reconstructed (Table 4). Neither pilot supports a claim that one prompt
or model is better than another, or any claim about unaided capability,
and none is made. Replies from hosted models can vary between calls even
at temperature 0.

One of the three models is undisclosed and may not remain available. Two
reports (C017 Arm 5 with
\texttt{meta-\allowbreak{}llama/\allowbreak{}llama-\allowbreak{}3.3-\allowbreak{}70b-\allowbreak{}instruct}
and C050 run 1) were overwritten in the repository history; byte copies
are now included in this package. The C043 and C044 model reports and
the C049 panel B runner remain unavailable in the version-controlled
evidence reviewed here. The repeated C050 panel shows that a single
reply per case can change between runs; C049 instead compares different
panels.

The package freshly replays the C017 cases and 15 retained solution
records. Its offline recount verifies source custody, extraction and
table arithmetic: Table 3 retains original acceptance flags, while Table
4 retains author annotations and the prior checker replay record. Fresh
C045--C050 checker execution still requires a repository checkout with
the cited commits. The C017 harness checkers share reference-solver
code. Lean formalizes only the Section 2.3 remark; the complete host is
tested, not mechanically formalized. This review did not execute the
pinned Lean compiler; historical Lean reports are retained without being
presented as fresh evidence.

This paper assigns no authorization and changes no jurisdiction. It
carries no operational authority or institutional standing.

\hypertarget{reproduction}{%
\subsection{6. Reproduction}\label{reproduction}}

From the package root with Python 3.12 or later, install the development
tools and run the offline checks (on Windows, substitute
\texttt{py -\allowbreak{}3.12} for \texttt{python}):

\begin{lstlisting}
python -m pip install -e '.[dev]'
python -m pytest code/tests
python -m ruff check code tools/study_capstone.py tools/render_paper.py
python -m mypy code
python code/replay_c017.py
python code/recount_pilots.py
python tools/study_capstone.py --output /tmp/paper010-capstone.json
python tools/render_paper.py
python verify_deposit.py
python tools/rebuild_release.py --check
\end{lstlisting}

Choose a writable capstone output path on Windows. The development extra
lists pytest 9 or later, ruff 0.14.13 and mypy 1.19.1; the current
receipt records the versions actually used. Runtime algebra and the
replay/recount scripts use only the standard library and packaged data.
They make no provider or network call.

The deposited renderer uses Pandoc and XeLaTeX to regenerate both
\texttt{paper.tex} and \texttt{paper.pdf}. To rebuild release custody
after an intentional edit, run
\texttt{python tools/\allowbreak{}rebuild\_\allowbreak{}release.py},
then both deposit checks again. Digest agreement establishes file
custody, not the truth of all manuscript claims.

With the exact toolchain in
\texttt{lean4/\allowbreak{}lean-\allowbreak{}toolchain} (Lean
4.34.0-rc1), run \texttt{lake build} in \texttt{lean4/\allowbreak{}},
then
\texttt{python tools/\allowbreak{}replay\_\allowbreak{}lean\_\allowbreak{}axiom\_\allowbreak{}census.py -\allowbreak{}-\allowbreak{}replay .}
from the package root. A static
\texttt{-\allowbreak{}-\allowbreak{}check .} only compares the generated
census source with the retained named report. The separate full-module
report includes generated structure declarations; it is not a second
mathematical theorem or a proof of the complete checker. The historical
C045--C050 checker replay needs the collection folder and cited Git
commits:
\texttt{python tools/\allowbreak{}replay\_\allowbreak{}formal\_\allowbreak{}math\_\allowbreak{}baselines.py -\allowbreak{}-\allowbreak{}check}.
That operation is distinct from this package's offline report recount.

\hypertarget{verification-for-this-revision}{%
\subsection{Verification for this
revision}\label{verification-for-this-revision}}

\texttt{docs/\allowbreak{}v01\_\allowbreak{}review\_\allowbreak{}2026\_\allowbreak{}10\_\allowbreak{}02.json}
records the reviewed repository baseline, commands, actual tool
versions, exit statuses and scope. Current tests, lint, type checking,
C017 replay, source/extraction/table recount, capstone, PDF rendering
and deposit custody checks were executed. The pinned Lean build and
fresh C045--C050 checker execution were not run in this environment, and
no model trial was made. The earlier
\texttt{docs/\allowbreak{}release\_\allowbreak{}validation\_\allowbreak{}*.json}
files and Lean reports retain their original dates and scope; they are
historical records. The current review record makes these limits
explicit rather than treating a retained report as execution.

\hypertarget{references}{%
\subsection{References}\label{references}}

\begin{enumerate}
\def\labelenumi{\arabic{enumi}.}
\tightlist
\item
  K. G. Wilson, ``Confinement of quarks,'' \emph{Physical Review D}
  \textbf{10}(8), 2445--2459 (1974).
  \href{https://doi.org/10.1103/PhysRevD.10.2445}{doi:10.1103/PhysRevD.10.2445}
\item
  J. B. Kogut, ``An introduction to lattice gauge theory and spin
  systems,'' \emph{Reviews of Modern Physics} \textbf{51}(4), 659--713
  (1979).
  \href{https://doi.org/10.1103/RevModPhys.51.659}{doi:10.1103/RevModPhys.51.659}
\item
  A. Hatcher, \emph{Algebraic Topology}. Cambridge University Press,
  2002, Chapter 1.
  \href{https://pi.math.cornell.edu/~hatcher/AT/ATpage.html}{Author's
  book page}
\item
  D. E. Knuth and P. B. Bendix, ``Simple word problems in universal
  algebras,'' in J. Leech (ed.), \emph{Computational Problems in
  Abstract Algebra}, pp.~263--297. Pergamon Press, 1970.
  \href{https://doi.org/10.1016/B978-0-08-012975-4.50028-X}{doi:10.1016/B978-0-08-012975-4.50028-X}
\item
  E. H. Bareiss, ``Sylvester's identity and multistep integer-preserving
  Gaussian elimination,'' \emph{Mathematics of Computation}
  \textbf{22}(103), 565--578 (1968).
  \href{https://doi.org/10.1090/S0025-5718-1968-0226829-0}{doi:10.1090/S0025-5718-1968-0226829-0}
\end{enumerate}

\end{document}
